% AERO2711 progress report -- LaTeX version.
% Compile with pdfLaTeX (Overleaf default) or tectonic. Figures are in figs/.
\documentclass[12pt,a4paper]{article}
\usepackage[utf8]{inputenc}
\usepackage[T1]{fontenc}
\usepackage{amsmath,amssymb}
\usepackage{newtxtext,newtxmath}
\usepackage[a4paper,left=3cm,right=2cm,top=2.54cm,bottom=2.54cm]{geometry}
\usepackage{setspace}
\onehalfspacing
\usepackage{graphicx}
\usepackage{float}
\usepackage{array,booktabs}
\usepackage{enumitem}
\usepackage[table]{xcolor}
\usepackage{textcomp}
\usepackage{listings}
\usepackage{cite}
\usepackage[skip=6pt]{caption}
\usepackage[hidelinks]{hyperref}
\lstset{basicstyle=\ttfamily\footnotesize,breaklines=true,columns=fullflexible}
\setlist{itemsep=2pt,topsep=4pt}
\setlength{\parskip}{4pt}
\newcommand{\ph}[1]{\colorbox{yellow}{#1}}

\begin{document}

\begin{titlepage}
  \centering
  \vspace*{3cm}
  {\large\color{gray} AERO2711 Research Project --- Progress Report\par}
  \vspace{1cm}
  {\LARGE\bfseries Visual Active Tracking of Manoeuvring Targets\\ under Occlusion\par}
  \vspace{0.6cm}
  {\large\itshape Map-constrained prediction and edge-aiming blind-phase control\\ for a drone-mounted camera\par}
  \vspace{3cm}
  \textbf{Student:} \ph{[Full name], SID [xxxxxxxxx]}\par
  \textbf{Supervisor:} \ph{[Title] Zi [Surname]}\par
  \textbf{Co-supervisor:} \ph{[name, or delete this line]}\par
  \vspace{1.5cm}
  \ph{[School / Faculty], [University]}\par
  September 2026\par
\end{titlepage}

\clearpage
\section*{Abstract}
\addcontentsline{toc}{section}{Abstract}
Visual active tracking (VAT) systems steer a camera to keep a person in view. When a wall or another obstacle hides the person, the tracker has to judge where the person will reappear, so that it can follow them again quickly. Current systems, including the recent OA-VAT framework, use a constant-velocity (CV) Kalman filter for this and are evaluated on targets that walk at near-constant speed. This project studies targets that speed up or slow down while they are hidden.

On a calibrated occlusion benchmark in Unreal Engine (an L-shaped corner), upgrading the motion model from CV to constant acceleration (CA) gave no significant gain. Most of the error came from how the filter is used during the occlusion. Partially occluded measurements corrupt its last estimate, and its blind-phase prediction drifts into walls. A reactive controller that steers toward this prediction then carries the error into the camera motion. This project tests a set of changes that use a map of the scene. It combines a \textbf{map-constrained Kalman filter (MC-KF)}, which repairs partially occluded measurements and projects the blind-phase prediction into reachable corridor space, with \textbf{edge-aiming blind-phase control}: while the target is hidden, the camera turns toward the edge of the obstacle instead of toward the filter's prediction and moves at full speed. Since no planner is used, in practice the camera slides along the wall toward the corner. Compared with the CV Kalman filter released with OA-VAT, run with the same ideal detector and controller on random target motion, the system reduces the re-acquisition error from 114 px to 25 px and is better on 31 of 33 paired seeds, with success rate 1.00 and no wall collisions.
The remaining lag while the target is hidden can be corrected either by the map projection or by a CA filter (12 px), but not by both, suggesting that, in this setting, the improvement comes from how the filter is used rather than from the motion model.

\clearpage
\tableofcontents
\clearpage
\section{Background and Motivation}
% --- Paragraph 1: the task and the problem ---
Small drones are increasingly used to follow a person automatically, for example for filming, inspection or search and rescue.
This task is called \emph{visual active tracking} (VAT): unlike a normal tracker, which only marks the person in a video, an active tracker also moves the camera to keep the person in view at a suitable distance \cite{r1,r2}.
Indoors and in cities, a common reason for failure is occlusion: a wall, a building or another obstacle hides the person, and the tracker must decide where to look and where to fly without seeing them.
The goal is to find the person again quickly once they reappear.

% --- Paragraph 2: how trackers handle this today ---
Most VAT systems bridge this gap with a Kalman filter that uses a constant-velocity (CV) model.
``Constant velocity'' is the filter's assumption, not a description of the person: the filter moves the last known position forward as if speed and direction stayed the same, and treats any difference as noise.

% --- Paragraph 3: the gap and what motivated this project ---
The recent OA-VAT system \cite{r3} uses such a filter, and the people in its tests walk at an almost constant speed, so this assumption is never really tested.
Real people, however, often speed up, stop or turn once they are out of sight; this is the usual way to escape a follower \cite{r4}.
Two observations made while reproducing OA-VAT (Section 6.1) pointed to this problem.
In a small test, gentle acceleration of the person lowered OA-VAT's tracking score by 14\%, and at corners its predictor did no better than using no predictor at all.
This project asks what goes wrong when a person changes speed while hidden, and how the tracker should handle it so that the person is found again quickly.

% --- Paragraph 4: first answer, the real cause, and the approach ---
The obvious answer is a better motion model, such as a constant-acceleration (CA) filter, which can follow a person who is speeding up.
Neither model matches the people here exactly, because they change speed at random moments, often while they are hidden.
In this project the CA filter helped only a little.
The larger errors come from three other sources.
First, while an obstacle hides part of the person, the detected box is cut off, so the filter is given a wrong position.
Second, while the person is fully hidden, the prediction can drift into a wall, where no person can be.
Third, if the camera turns toward this wrong prediction, the error goes straight into where the camera looks and flies.
Since the layout of walls and corridors is usually known or can be mapped, this project uses a map of the surroundings to correct the filter and to guide the camera, in a way that works with any motion model.
The method is developed and tested on an L-shaped corner, the simplest layout that produces a repeatable occlusion.

\section{Literature Survey}
\subsection{Visual active tracking}
Luo et al. \cite{r1,r2} trained a single neural network with reinforcement learning to turn camera images directly into camera movements, using virtual worlds built with UnrealCV \cite{r5}.
Later work made this approach more robust: AD-VAT \cite{r6} trains the tracker against a target that learns to escape, \cite{r7} adds people who look similar to the target, RSPT \cite{r8} rebuilds the surroundings and predicts the target's path so that the tracker works in new scenes, and \cite{r9} uses large pre-trained vision models with reinforcement learning from recorded data.
In all of these methods, the handling of occlusion is hidden inside the network's memory, so when tracking fails it is hard to tell whether the prediction or the control was wrong.
OA-VAT \cite{r3} is instead built from separate modules.
It finds and identifies the person with YOLOE \cite{r10}, DINOv3 features \cite{r11} and ORTrack \cite{r12}.
When the person is lost, it stops the camera, predicts where the person has gone with a Kalman filter that weighs each measurement by the detector's confidence, searches there, and hands over to a diffusion-policy planner \cite{r13} if the person is still not found.
This project reuses its Kalman filter, evaluation protocol and simulator (Section 5.1), but not its tracking pipeline.

\subsection{Motion models and occlusion in tracking}
The Kalman filter \cite{r14} is optimal when the motion model matches the target. For targets that do not follow a fixed model, the tracking literature offers the Singer acceleration model \cite{r15}, a family of dynamic models surveyed by Li and Jilkov \cite{r16}, and the Interacting Multiple Model (IMM) estimator \cite{r17}, the standard tool in radar and aerospace tracking \cite{r18}. All of these assume a steady stream of measurements. During a long occlusion there are no updates, and every model reduces to open-loop extrapolation from the estimate at the moment the target disappeared.

In computer vision, the SORT family \cite{r19,r20,r21} runs a CV Kalman filter on the bounding box to bridge missed detections. OC-SORT \cite{r22} shows that a filter coasting through occlusion builds up error, mainly through its noisy velocity estimate, and corrects the filter \emph{after} the target is found again. ORTrack \cite{r12} makes the appearance model robust to partial occlusion.
A related, well-known issue is that a box cut off by an obstacle has its centre shifted toward the visible side, just before the target disappears.
Trackers usually handle this by trusting such measurements less, as OA-VAT and StrongSORT \cite{strongsort} do by scaling the measurement noise with detector confidence, or by training detectors to predict the full extent of a partly hidden object \cite{amodal}.
When the layout of the scene is known, a simple complementary option is to use it directly; this project explores rebuilding the box centre from its visible edge, with the map indicating which side is cut off.

\subsection{Scene-aware prediction and constrained estimation}
Human trajectory prediction has moved to learned models that take the surroundings into account \cite{r23}, such as Social LSTM \cite{r24} and Trajectron++ \cite{r25}.
Kitani et al. \cite{r26} showed that knowing where a person \emph{can} walk strongly narrows down long-range forecasts.
These predictors are mostly evaluated offline on recorded trajectories, rather than inside a tracking loop where the camera itself moves.
In estimation theory, known physical limits can be added to a Kalman filter; this is called \emph{state-constrained Kalman filtering} \cite{r27}.
The simplest method, estimate projection, moves an estimate that falls outside the allowed region to the nearest point inside it.
The MC-KF in this project applies this idea, with the free space in a map of the surroundings (here, the corridor) as the allowed region, after converting the image position to a floor position.

\subsection{Aerial target following among obstacles}
Model-based aerial trackers plan drone trajectories with explicit collision and visibility terms \cite{r28,r29,r30}. Fast-Tracker \cite{r31} and Elastic Tracker \cite{r32} track agile targets in clutter using a map of the environment \cite{r33} and a kinematic prediction of the target. They are designed to \emph{prevent} occlusion, and they assume a metric map and a 3D target state. Collision safety is often enforced with control barrier functions \cite{r34}, which change a nominal command as little as possible. The safety layer used here follows the same idea at velocity level, and never changes the heading.

\subsection{Summary}
In short, most trackers bridge occlusion with a constant-velocity filter, and common VAT benchmarks rarely test targets that change speed while hidden.
Known scene layouts are widely used in trajectory prediction and aerial planning, but less often inside the estimator of an active tracker.
This project therefore tests, in a small simulation study, whether a map of the surroundings can improve re-acquisition after occlusion when the target changes speed.

\section{Problem Statement}
\textbf{In a simulated drone tracking system, how can the processing of measurements, the prediction and the camera's motion during occlusion be arranged, using a map of the surroundings, so that a person who changes speed while hidden behind an obstacle is found again quickly and without collisions?}

The design target is a re-acquisition error below about 30 px, with success rate 1.00 and zero wall contacts, for targets that move along random routes at random speeds.

\subsection{Work to date}
\begin{enumerate}[label=\textbf{C\arabic*.}]
  \item \textbf{Map-constrained Kalman filter (MC-KF).} A standard box-space Kalman filter with two map-based steps: repairing partially occluded measurements, and projecting the blind-phase prediction into reachable corridor space. It works with either a CV or a CA motion model.
  \item \textbf{Edge-aiming blind-phase control.} While the target is hidden, the camera turns toward the edge of the obstacle, where the target will reappear, instead of toward the filter's prediction, and moves at full speed. A velocity-level safety layer keeps it clear of walls. This is a simple rule used in place of a planner; in practice the camera slides along the wall toward the corner.
  \item \textbf{Steering during occlusion.} When the camera kept turning toward the filter's prediction without the map projection, the tracker lost the target in 6 of 40 episodes. With the projection, following the prediction no longer lost the target, but the re-acquisition error was about 100 px larger than when aiming at the obstacle edge. In this setting, aiming at the edge gave the most accurate re-acquisition.
  \item \textbf{Benchmark evidence.} A calibrated occlusion benchmark (an L-shaped corner), a re-acquisition metric and a full set of ablations. They show that CV $\rightarrow$ CA gives no significant gain on the baseline, and that inside the full system the CA acceleration term and the map projection correct the same lag, so only one of them should be used.
\end{enumerate}
\section{Project Objectives}
$\checkmark$ = complete, $\boxminus$ = in progress, $\square$ = remaining.

\begin{enumerate}
  \item $\checkmark$ Reproduce OA-VAT \cite{r3} in simulation as the starting point for understanding current methods.
  \item $\checkmark$ Build a calibrated L-corner test with engine-verified occlusion timing and a closed-loop test program that runs 40 paired seeds.
  \item $\checkmark$ Test whether a better motion model helps (no predictor / CV / CA, with the original and the calibrated measurement noise).
  \item $\checkmark$ Design the MC-KF and edge-aiming blind-phase control, and test each part by removing it.
  \item $\boxminus$ Move to self-built mazes in UnrealZoo, with a general map module that computes the free space and occluding edges from the maze layout instead of the hand-coded L-corner values. \emph{Test:} no predictor / CV / CA with random target motion, 40 seeds per maze; report SR, AR, re-acquisition error in pixels and time to re-acquire.
  \item $\square$ Add a planner for the camera's motion while the target is hidden, replacing the edge-aiming rule, to form a complete tracking system. First maze runs show that the simple rule is not enough there: the camera reaches the next turn too fast to follow the target. \emph{Test:} edge-aiming rule vs planner on the same mazes and seeds, with the same metrics.
  \item $\square$ Handle junctions, where the target may leave by more than one exit: keep several predictions and let the planner decide which exit to check first.
  \item $\square$ Replace the oracle detector with the OA-VAT perception stack; re-calibrate the measurement noise and test rejecting detections that do not fit the prediction.
  \item $\square$ Study braking targets (short occlusions, overshoot); add an IMM baseline and repeat the main comparisons over three seed sets with confidence intervals.
  \item $\square$ Stretch: indoor test on a Tello drone.
\end{enumerate}
\section{Methodology}
\subsection{Simulation setup}
The benchmark runs in the FlexibleRoom scene of gym-unrealcv \cite{r5}. Five scripted blocks form an L-shaped corridor, and each wall position is checked against the engine. No planner is used: the camera is steered by a simple reactive controller, so any difference in performance comes from the estimation and blind-phase logic (Figure 1). Table 1 lists the main settings.

\begin{figure}[H]
  \centering
  \includegraphics[width=15.5cm]{figs/fig_arch.png}
  \caption{One control step of the tracking system. Orange boxes are the parts added in this project: measurement repair and map projection form the MC-KF, and ``Steering: hidden'' is the edge-aiming rule used while the target is hidden.}
\end{figure}

\begin{table}[htbp]
  \centering
  \small
  \caption{Frozen benchmark configuration.}
  \renewcommand{\arraystretch}{1.15}
  \begin{tabular}{>{\raggedright\arraybackslash}p{0.247\linewidth}>{\raggedright\arraybackslash}p{0.295\linewidth}>{\raggedright\arraybackslash}p{0.368\linewidth}}
  \toprule
  \textbf{Item} & \textbf{Value} & \textbf{Reason} \\
  \midrule
  Corridor & W = 150 cm, walls 350 cm & Narrow indoor corridor \\
  Camera & 640$\times$480, hFOV 82.6\textdegree{}, pitch $-$15\textdegree{} & Tello spec; target centred at 250 cm \\
  Camera limits & v $\le$ 99 cm/s, |$\omega$| $\le$ 30\textdegree{}/s, 20 Hz & Speed and rate from OA-VAT; turn-rate limit set here \\
  Target & Person, 168 cm, speed 40–99 cm/s & Never faster than the camera \\
  Following distance & 250 cm & Optimum of the OA-VAT score \\
  Wall margin & camera $\ge$ 30 cm from walls & Near-clip limit 15 cm (measured) plus body radius and localisation slack \\
  \bottomrule
  \end{tabular}
\end{table}

\textbf{What is taken from OA-VAT.} Only the elements in Table 2 are inherited. Everything else is designed in this project: the track, camera setup, detector, controller, target behaviours, re-acquisition metric and the proposed architecture.

\begin{table}[htbp]
  \centering
  \small
  \caption{Elements inherited from OA-VAT.}
  \renewcommand{\arraystretch}{1.15}
  \begin{tabular}{>{\raggedright\arraybackslash}p{0.223\linewidth}>{\raggedright\arraybackslash}p{0.416\linewidth}>{\raggedright\arraybackslash}p{0.271\linewidth}}
  \toprule
  \textbf{Element} & \textbf{Inherited content} & \textbf{Difference in this project} \\
  \midrule
  Kalman filter & 8-D CV filter from the OA-VAT code (structure from DeepSORT \cite{r20}); noise weights 1/20 (position), 1/40 (velocity) & Used unchanged as the baseline; the CA filter extends it \\
  Evaluation protocol & Per-step score r = 1 $-$ |$\rho$ $-$ 250|/750 $-$ |$\theta$|/45; visible if $\rho$ $\le$ 750 cm and |$\theta$| $\le$ 45\textdegree{}; 500 steps; AR / EL / SR & Failure after 50 invisible steps (OA-VAT: after more than 50); r clipped to [$-$1, 1] before summing \\
  Control timing & 50 ms period (20 Hz), speed limit 99 cm/s, following distance 250 cm & None \\
  Simulator & gym-unrealcv FlexibleRoom scene and person model & Windows fixes; walls built from scripted blocks \\
  \bottomrule
  \end{tabular}
\end{table}

\textbf{Scope of the comparison.} Perception is simplified in the same way as control: detection uses the engine’s object mask instead of OA-VAT’s YOLOE, DINOv3 and ORTrack stack, and steps run in fixed 50 ms increments with no latency. OA-VAT’s own occlusion procedure is therefore not reproduced here. That procedure stops the camera, searches around the prediction with its detector for at most 15 frames, and then falls back to the planner. All results in this report compare configurations under the same detector and controller.


\subsection{Occlusion geometry}
A simple model is used to choose the following distance. The target turns the corner at a distance c from the inner corner, and the camera follows along the same path, a path length d behind; the target is treated as a point. When the target has walked a distance s past the corner, the wall blocks the view if $s^2 - d\,s + c\,d < 0$. This happens only if

\begin{equation}
  d > 4c
\end{equation}
and the target stays hidden over a path length

\begin{equation}
  L_{\mathrm{occ}} = \sqrt{d\,(d-4c)}
\end{equation}
With d = 250 cm and c = 40 cm this gives 150 cm (about 38 steps at 80 cm/s), which matches the engine's object mask. Each passage has a \emph{fade-out} phase, in which the wall gradually clips the box, followed by a \emph{fully blind} phase (Figure 2).
The model holds only when the camera follows the same path; a camera closer to the inner wall is blocked earlier, and one near the outer wall may not be blocked at all. It is therefore used only to choose d = 250 cm and c in [40, 60] cm; whether an episode contains an occlusion is measured from the object mask (Section 5.3).

\begin{figure}[H]
  \centering
  \includegraphics[width=14.5cm]{figs/fig_occlusion.png}
  \caption{Visible target pixels measured in the engine as the target walks around the corner (d = 250 cm, c = 40 cm, camera on the target's path, 80 cm/s). Shaded: the fade-out phase, in which the wall gradually cuts off the box, and the fully hidden phase.}
\end{figure}

\subsection{Target behaviour}
Each seed generates a random route through the corner and a random speed profile: every 2–4 s the target picks a new speed between 40 and 99 cm/s and changes to it with an acceleration of at most 80 cm/s\textsuperscript{2}. Nothing is tied to the corner, so the target may or may not change speed while it is hidden.

Because the route and speed are random, not every episode contains an occlusion, so the episodes fall into two groups, and results are reported for each:
\begin{itemize}
  \item \textbf{With occlusion} (33 of 40 seeds): the target is completely hidden for at least 5 frames. Only these episodes have a re-acquisition error.
  \item \textbf{Without occlusion} (7 of 40 seeds): the target is never completely hidden. These episodes check that the system does not harm ordinary tracking, using AR, SR and wall contacts.
\end{itemize}
The split is taken from the full system with CV; the baseline has one more occluded seed (34), because its camera follows a slightly different path.

\subsection{Baseline filter and controller}
\textbf{State.} The filter tracks the target's bounding box in the image. Its state is
\begin{equation}
  \mathbf{x} = [\,c_x,\ c_y,\ a,\ h,\ \dot c_x,\ \dot c_y,\ \dot a,\ \dot h\,]^{\top}
\end{equation}
and the measurement is the detected box $\mathbf{z} = [\,c_x,\ c_y,\ a,\ h\,]^{\top}$. Table 3 lists every entry, including the two added by the CA variant. One step is one frame (50 ms). The noise is scaled by the box height $h$, because a larger (closer) box also moves more in pixels.

\begin{table}[H]
  \centering
  \small
  \caption{Filter state entries and noise settings (standard deviations; the CV state has the first eight entries, the CA state all ten).}
  \renewcommand{\arraystretch}{1.15}
  \begin{tabular}{llllll}
  \toprule
  \textbf{Symbol} & \textbf{Meaning} & \textbf{Unit} & \textbf{Process noise} & \textbf{Measured (noise)} & \textbf{Used in} \\
  \midrule
  $c_x,\ c_y$ & box centre & px & $h/20$ & yes ($h/20$) & CV, CA \\
  $a$ & width-to-height ratio & -- & 0.01 & yes (0.1) & CV, CA \\
  $h$ & box height & px & $h/20$ & yes ($h/20$) & CV, CA \\
  $\dot c_x,\ \dot c_y$ & velocity of the centre & px/frame & $h/40$ & no & CV, CA \\
  $\dot a$ & change of the ratio & 1/frame & $10^{-5}$ & no & CV, CA \\
  $\dot h$ & change of the height & px/frame & $h/40$ & no & CV, CA \\
  $\ddot c_x,\ \ddot c_y$ & acceleration of the centre & px/frame\textsuperscript{2} & $h/800$ & no & CA only \\
  \bottomrule
  \end{tabular}
\end{table}

\textbf{Model.} The CV model assumes that the four velocities stay constant:
\begin{equation}
  \mathbf{x}_{k+1} = F\mathbf{x}_k + \mathbf{w}_k,\quad \mathbf{z}_k = H\mathbf{x}_k + \mathbf{v}_k,\quad F=\begin{bmatrix} I_4 & I_4\\ 0 & I_4\end{bmatrix},\quad H=\begin{bmatrix} I_4 & 0\end{bmatrix}
\end{equation}
so each position entry is advanced by its velocity, and only the box itself is measured.

\textbf{Baseline.} The baseline is this filter as released with OA-VAT (Table 2), with its original noise settings: it predicts every frame, updates whenever the target is detected, and coasts until the target reappears.

\textbf{CA variant.} The CA state adds the two accelerations of the box centre:
\begin{equation}
  \mathbf{x}_{\mathrm{CA}} = [\,c_x,\ c_y,\ a,\ h,\ \dot c_x,\ \dot c_y,\ \dot a,\ \dot h,\ \ddot c_x,\ \ddot c_y\,]^{\top}
\end{equation}
The box centre is advanced with both its velocity and its acceleration, while $a$ and $h$ keep the CV model. The accelerations are not measured; they are learned from how the box centre moves, and are limited to the pedestrian bound of 80 cm/s\textsuperscript{2}. Converted to pixels with the box height (a 168 cm person appears $h$ px tall) and one frame of 0.05 s, this gives
\begin{equation}
  |\ddot c_x|,\ |\ddot c_y| \le 80 \cdot \frac{h}{168} \cdot 0.05^2 \approx 0.0012\,h \ \ \text{px per frame}^2.
\end{equation}

\textbf{Controller.} When the target is visible, the camera turns toward the estimated bearing $\theta$ and keeps the estimated range $\hat\rho$ at 250 cm:
\begin{equation}
  \theta = \arctan\frac{c_x-320}{f_x},\qquad \hat{\rho} = \frac{f_x\cdot 168}{h}
\end{equation}
\begin{equation}
  \omega = \operatorname{clip}(3\theta,\ \pm 30^\circ/\mathrm{s}),\qquad v = \operatorname{clip}\big(2(\hat{\rho}-250),\ 0,\ 99\ \mathrm{cm/s}\big)
\end{equation}
Here 320 px is the image centre column, $f_x = 364.3$ px is the focal length, and 168 cm is the person's height, so $\hat\rho$ is the pinhole-camera range estimate. $\omega$ is the turn rate and $v$ the forward speed; the gains 3 and 2 were chosen by hand. In the baseline the same controller is also driven by the filter's prediction while the target is hidden, so the camera keeps steering toward the extrapolated box.

\subsection{Map-constrained Kalman filter (MC-KF)}
The MC-KF keeps the standard predict and update equations and adds two steps that use the corridor map.

\textbf{Measurement repair (fade-out frames).} When the target’s pixel count drops below 50\% of its peak over the last 8 healthy frames, the box is being clipped by the wall and its centre is biased. The corner apex is projected into the image to decide which side is clipped. The centre is rebuilt from the clean edge using the box width remembered from healthy frames: cx = u\textsubscript{clean} $\pm$ w\textsubscript{good}/2. Only this repaired cx is used in the update. Figure 3 shows a real fade-out frame.

\begin{figure}[H]
  \centering
  \includegraphics[width=10cm]{figs/fig_repair.png}
  \caption{Measurement repair on a real fade-out frame (camera view, full system with CV). The wall on the left cuts off part of the target, so the detected box (red) and its centre are shifted to the right. The occluding edge (orange) is computed from the map; the box is rebuilt from its visible right edge with the width remembered from healthy frames (green). In this frame the error to the true centre (blue) falls from 46 px to 24 px.}
\end{figure}

\textbf{Reachable-set projection (blind frames).} After each prediction, the predicted (cx, h) is back-projected to a floor point using the camera pose $\psi$:

\begin{equation}
  \mathbf{p} = \mathbf{p}_{\mathrm{cam}} + \hat{\rho}\,\big[\cos(\psi+\theta),\ \sin(\psi+\theta)\big]^{\top}
\end{equation}
If p lies outside the corridor free space, it is moved to the nearest reachable point and projected back into (cx, h). The covariance is not changed, so the filter stays cautious and trusts the first new measurement. This is estimate projection in the sense of \cite{r27}. Two earlier versions, pushing the point out of the nearest wall and ray-marching through walls, failed and were discarded. Figure 4 shows the first blind frame of the same episode.

\begin{figure}[H]
  \centering
  \includegraphics[width=10cm]{figs/fig_clamp.png}
  \caption{Map projection on the first blind frame of the same episode (top view). The filter prediction (red cross), placed along the predicted bearing and range, lies inside the wall; it is moved to the nearest point of the free space (green). Its distance to the true target position (blue) falls from 106 cm to 70 cm.}
\end{figure}

\subsection{Edge-aiming blind-phase control and safety}
While the target is fading or hidden, yaw aims at the image column of the occluding corner edge, where the target must reappear, and speed is fixed at 99 cm/s. The filter's estimate (its \emph{belief}) is no longer used for steering. This is a simple rule chosen because the test bench uses no planner. Combined with the safety layer below, its visible effect is that the camera slides along the wall toward the corner; a planner could choose better viewpoints, for example moving out to see around the corner earlier. It keeps running in the background, is scored when the target reappears, and then continues with the new measurement. A safety layer limits only the wall-approach component of the velocity, based on the distance d to the nearest wall:

\begin{equation}
  v_{\mathrm{in}} \le \frac{\max(0,\ d-30\,\mathrm{cm})}{0.6\,\mathrm{s}}
\end{equation}
The safety layer never changes heading. A hard veto with wall sliding backs it up.

\subsection{Evaluation}
On a single corner the OA-VAT metrics (accumulated reward AR, episode length EL, success rate SR) change very little between methods, because every method eventually finds the target again. The main metric is therefore the \textbf{re-acquisition error}. When the target is detected again after at least 5 blind frames, it is the horizontal pixel distance between the filter’s predicted box centre and the centre of the newly detected box. All errors are reported in pixels, because the distance to the target at reappearance varies and a conversion to centimetres would be unreliable. Every configuration runs on the same 40 seeds (2000–2039) with identical target routes. Results are reported as paired differences, per-seed win counts and bootstrap 95\% confidence intervals.

\section{Progress Report and Revised Gantt Chart}
\subsection{Reproduction and motivating observations}
OA-VAT was reproduced on Windows with all three modules, including training the diffusion planner on 1,464 self-collected episodes. This required several fixes to the released code. One of them was a missing dependency that had silently disabled the confidence-aware smoothing. On UrbanCity the reproduction reached SR 0.95 and EL 487 (paper: 0.95 / 486), and on ParkingLot the planner raised SR from 0.70 to 0.90 (paper: 0.93). AR was 80–115 points lower than reported (270 vs 385 on UrbanCity, 314 vs 392 on ParkingLot), and the matching SR values were obtained before the missing dependency was found. The reproduction is therefore close but not exact. It served to understand the method rather than to benchmark it.

Two observations from this stage motivated the project. On a four-corner course, target acceleration capped at the camera’s top speed reduced AR from 291 to 249 (10 episodes, one seed). In a ParkingLot variant where the target manoeuvres at corners, the confidence-aware predictor reached SR 0.85 against 0.90 without a predictor (20 episodes). Both samples are small, so these results are treated as motivation, not as evidence.

\subsection{The motion-model hypothesis}
The original plan was to replace CV with a stronger motion model. This was tested on the calibrated L-corner with 40 paired seeds of random target motion; 34 of them contain a full blind phase. With the inherited noise settings, CA and CV gave the same error (difference $-$0.7 px, 95\% CI [$-$5.8, 4.7]). With the measurement noise reduced to match the near-noiseless detector, CA was 2.7 px better than CV (98.5 vs 101.3 px), but the difference is not significant (CI [$-$6.3, 1.8]). Calibrating the noise itself helped a little (CV: 114.1 $\rightarrow$ 101.3 px). In most episodes the target did not change speed while hidden: the median speed change during the blind phase was 4.5 cm/s, and only 7 of 33 episodes changed by more than 15 cm/s. The error therefore comes mainly from the estimate at the moment the target disappears, and a better motion model can do little about it.

\subsection{Measurement repair and steering during occlusion}
The next step was to clean the measurements. In early single-seed tests, a more accurate repaired estimate made the camera turn deep behind the wall and miss the target reappearing at the corner edge.
A 40-seed control supports this observation: when the camera followed the filter's prediction and no map projection was used, adding measurement repair lowered the error on the surviving episodes by 75 px, but the tracker lost the target in 6 of 40 episodes (SR 1.00 $\rightarrow$ 0.85, AR 455 $\rightarrow$ 428).
With the map projection added, following the prediction no longer lost the target (SR 1.00), but the re-acquisition error was 101 px larger than with edge aiming (CI [82, 119]; 32 of 33 seeds worse).
Aiming at the obstacle edge therefore gave the most accurate re-acquisition in this setting, and together with measurement repair it gave the largest improvement in this study. Figure 5 shows the sequence.

\begin{figure}[htbp]
  \centering
  \includegraphics[width=14.5cm]{figs/fig_progress.png}
  \caption{Re-acquisition error at each development stage (random target motion, 40 paired seeds). Stages 2–5 use the calibrated measurement noise.}
\end{figure}

Measured with CV, the reduction from 114 px to 24 px breaks down as follows: noise calibration 13 px, velocity-cap safety layer 12 px, measurement repair with edge aiming 61 px, and fixed blind speed with the reachable-set projection 5 px. With the original noise settings the full system reaches 25 px, so it does not depend on the noise calibration.

\subsection{Current results}
\begin{table}[htbp]
  \centering
  \small
  \caption{Current results with random target motion (30 cm wall margin, 40 seeds). With occlusion: 33 seeds (34 for the baseline); without occlusion: 7 seeds. SR is 1.00 in every case.}
  \renewcommand{\arraystretch}{1.15}
  \begin{tabular}{>{\raggedright\arraybackslash}p{0.36\linewidth}>{\raggedright\arraybackslash}p{0.13\linewidth}>{\raggedright\arraybackslash}p{0.09\linewidth}>{\raggedright\arraybackslash}p{0.09\linewidth}>{\raggedright\arraybackslash}p{0.09\linewidth}>{\raggedright\arraybackslash}p{0.09\linewidth}}
  \toprule
   & \multicolumn{3}{l}{\textbf{With occlusion}} & \multicolumn{2}{l}{\textbf{Without occlusion}} \\
  \textbf{Configuration} & \textbf{Re-acq. error (px)} & \textbf{AR} & \textbf{Wall hits} & \textbf{AR} & \textbf{Wall hits} \\
  \midrule
  Baseline: OA-VAT CV-KF (no map, no safety layer) & 114.1 & 450.9 & 53.3 & 453.4 & 81.9 \\
  Full system, no predictor & 93.1 & 457.8 & 0 & 459.5 & 0 \\
  \textbf{Full system, CV with map projection} & \textbf{24.6} & 457.8 & 0 & 459.6 & 0 \\
  \textbf{Full system, CA without map projection} & \textbf{11.6} & 457.8 & 0 & 459.5 & 0 \\
  Full system, CA with map projection & 38.2 & 457.8 & 0 & 459.5 & 0 \\
  \bottomrule
  \end{tabular}
\end{table}

Compared with the baseline, the full system with CV and the map projection lowers the error by 92 px (CI [$-$114, $-$71]; better on 31 of 33 seeds) and removes all wall contacts. The Kalman filter is still needed: with the same system but no predictor, the error is 93 px. AR and SR barely change, because on one corner every method eventually finds the target.

\textbf{Two ways to correct the lag.} While the target is hidden, a CV filter lags behind it. Two parts of the system can correct this lag: the map projection, which pulls the prediction back into the free space, and the acceleration state of a CA filter, which learns the motion trend before the target disappears. Each works on its own. CV with the projection reaches 24.6 px, and CA without the projection reaches 11.6 px, which is the lowest error in this study (13 px lower than CV with the projection; CI [$-$20, $-$6]; better on 23 of 33 seeds). Using both at once over-corrects: CA with the projection is 13.6 px worse than CV with the projection (CI [11, 16]; CA better on only 1 of 33 seeds). Setting the CA acceleration to zero during the blind phase removes most of this penalty (+3.2 px). For CV, the projection itself has no measurable effect under random motion ($-$0.6 px, CI [$-$12, 11]); it helped more in earlier tests where the target sprinted at the corner. \textbf{In this setting, the improvement comes from measurement repair and edge aiming, while the lag can be handled either by the map or by the motion model, but not by both.}

\textbf{Ablations.} Table 5 removes one component at a time from the full system with CV. Measurement repair has the largest effect in these tests: without it the error rises to about 340 px, worse than using no predictor, because the cut-off box is fed straight into the filter. Aiming at the occluding edge and flying at fixed speed are both needed. Edge aiming lowers AR slightly, since the score rewards keeping the target centred, but improves re-acquisition by 40 px. The velocity-cap safety layer is a fixed requirement and is kept in every configuration; removing it in the ablation only measures its cost: without it the camera hits a wall 59 times per episode, while the re-acquisition error is about 6 px lower. Zero collisions therefore cost about 6 px of accuracy.

\begin{table}[htbp]
  \centering
  \small
  \caption{Leave-one-out ablation, random target motion, 40 paired seeds (n.s. = confidence interval includes zero).}
  \renewcommand{\arraystretch}{1.15}
  \begin{tabular}{>{\raggedright\arraybackslash}p{0.40\linewidth}>{\raggedright\arraybackslash}p{0.25\linewidth}>{\raggedright\arraybackslash}p{0.25\linewidth}}
  \toprule
  \textbf{Removed from full system (CV)} & \textbf{$\Delta$ re-acq. error (px, 95\% CI)} & \textbf{Other effect} \\
  \midrule
  Measurement repair & +314 [+293, +336] & — \\
  Fixed blind speed & +44 [+36, +51] & AR $-$18 \\
  Edge aiming & +40 [+31, +51] & AR +8 \\
  Map projection & $-$1 [$-$12, +11] (n.s.) & — \\
  Velocity-cap safety layer (required; removed only to measure its cost) & $-$6 [$-$9, $-$4] & 59 wall contacts per episode \\
  \bottomrule
  \end{tabular}
\end{table}

\textbf{Limitations.} The results use an oracle detector and a single corner with a hand-coded map. They show the benefit of the system over the OA-VAT filter under identical conditions, not a head-to-head comparison with the full OA-VAT system, which would require both pipelines to run on the same real detector. With random motion, the target changed speed noticeably while hidden in only 7 of 33 episodes, so the results say more about occlusion in general than about speed changes in particular. Repeating a configuration gives small run-to-run differences, so only paired differences with confidence intervals are interpreted.

\subsection{Revised Gantt chart}
The original plan had five stages: a manoeuvring target, evidence that prediction degrades, CA/IMM baselines, a learned predictor, and a final comparison. The first three stages are done except for IMM. The results showed that the gains come from the map and from separating estimation and control, not from the motion model. The MC-KF architecture therefore became the main focus of the work, and the learned predictor became a stretch goal (Figure 6).

\begin{figure}[htbp]
  \centering
  \includegraphics[width=15.5cm]{figs/fig_gantt.png}
  \caption{Revised Gantt chart (green: complete; orange: in progress; blue: planned; hatched: stretch).}
\end{figure}

\section{Planned Work}
The final report will add the following sections:

\subsection{General map representation and multi-corner tracks}
\subsection{Real detector in the loop and innovation gating}
\subsection{Braking targets and blind-phase speed}
\subsection{Statistical validation and IMM baseline}
\subsection{Discussion and conclusion}
\clearpage
\begin{thebibliography}{99}
\addcontentsline{toc}{section}{References}
\bibitem{r1} W. Luo, P. Sun, F. Zhong, W. Liu, T. Zhang, and Y. Wang, ``End-to-end active object tracking via reinforcement learning,'' in \emph{Proc. Int. Conf. Machine Learning (ICML)}, 2018, pp. 3286–3295.
\bibitem{r2} W. Luo, P. Sun, F. Zhong, W. Liu, T. Zhang, and Y. Wang, ``End-to-end active object tracking and its real-world deployment via reinforcement learning,'' \emph{IEEE Trans. Pattern Anal. Mach. Intell.}, vol. 42, no. 6, pp. 1317–1332, 2020.
\bibitem{r3} H. Sun, K. Zhou, H. Gao, S. Zhang, J. Hu, X. Wen, Q. Ye, and M. Tan, ``Instance-level visual active tracking with occlusion-aware planning,'' arXiv:2604.21453, 2026.
\bibitem{r4} T. H. Chung, G. A. Hollinger, and V. Isler, ``Search and pursuit-evasion in mobile robotics: A survey,'' \emph{Autonomous Robots}, vol. 31, no. 4, pp. 299–316, 2011.
\bibitem{r5} W. Qiu, F. Zhong, Y. Zhang, S. Qiao, Z. Xiao, T. S. Kim, Y. Wang, and A. Yuille, ``UnrealCV: Virtual worlds for computer vision,'' in \emph{Proc. ACM Int. Conf. Multimedia}, 2017, pp. 1221–1224.
\bibitem{r6} F. Zhong, P. Sun, W. Luo, T. Yan, and Y. Wang, ``AD-VAT: An asymmetric dueling mechanism for learning visual active tracking,'' in \emph{Proc. Int. Conf. Learning Representations (ICLR)}, 2019.
\bibitem{r7} F. Zhong, P. Sun, W. Luo, T. Yan, and Y. Wang, ``Towards distraction-robust active visual tracking,'' in \emph{Proc. Int. Conf. Machine Learning (ICML)}, 2021, pp. 12782–12792.
\bibitem{r8} F. Zhong, X. Bi, Y. Zhang, W. Zhang, and Y. Wang, ``RSPT: Reconstruct surroundings and predict trajectory for generalizable active object tracking,'' in \emph{Proc. AAAI Conf. Artificial Intelligence}, 2023.
\bibitem{r9} F. Zhong, K. Wu, H. Ci, C. Wang, and H. Chen, ``Empowering embodied visual tracking with visual foundation models and offline RL,'' in \emph{Proc. European Conf. Computer Vision (ECCV)}, 2024.
\bibitem{r10} A. Wang, L. Liu, H. Chen, Z. Lin, J. Han, and G. Ding, ``YOLOE: Real-time seeing anything,'' in \emph{Proc. IEEE/CVF Int. Conf. Computer Vision (ICCV)}, 2025.
\bibitem{r11} O. Siméoni \emph{et al.}, ``DINOv3,'' arXiv:2508.10104, 2025.
\bibitem{r12} Y. Wu, X. Wang, X. Yang, M. Liu, D. Zeng, H. Ye, and S. Li, ``Learning occlusion-robust vision transformers for real-time UAV tracking,'' in \emph{Proc. IEEE/CVF Conf. Computer Vision and Pattern Recognition (CVPR)}, 2025, pp. 17103–17113.
\bibitem{r13} C. Chi, S. Feng, Y. Du, Z. Xu, E. Cousineau, B. Burchfiel, and S. Song, ``Diffusion policy: Visuomotor policy learning via action diffusion,'' in \emph{Proc. Robotics: Science and Systems (RSS)}, 2023.
\bibitem{r14} R. E. Kalman, ``A new approach to linear filtering and prediction problems,'' \emph{J. Basic Engineering}, vol. 82, no. 1, pp. 35–45, 1960.
\bibitem{r15} R. A. Singer, ``Estimating optimal tracking filter performance for manned maneuvering targets,'' \emph{IEEE Trans. Aerospace and Electronic Systems}, vol. AES-6, no. 4, pp. 473–483, 1970.
\bibitem{r16} X. R. Li and V. P. Jilkov, ``Survey of maneuvering target tracking. Part I: Dynamic models,'' \emph{IEEE Trans. Aerospace and Electronic Systems}, vol. 39, no. 4, pp. 1333–1364, 2003.
\bibitem{r17} H. A. P. Blom and Y. Bar-Shalom, ``The interacting multiple model algorithm for systems with Markovian switching coefficients,'' \emph{IEEE Trans. Automatic Control}, vol. 33, no. 8, pp. 780–783, 1988.
\bibitem{r18} Y. Bar-Shalom, X. R. Li, and T. Kirubarajan, \emph{Estimation with Applications to Tracking and Navigation}. New York: Wiley, 2001.
\bibitem{r19} A. Bewley, Z. Ge, L. Ott, F. Ramos, and B. Upcroft, ``Simple online and realtime tracking,'' in \emph{Proc. IEEE Int. Conf. Image Processing (ICIP)}, 2016, pp. 3464–3468.
\bibitem{r20} N. Wojke, A. Bewley, and D. Paulus, ``Simple online and realtime tracking with a deep association metric,'' in \emph{Proc. IEEE Int. Conf. Image Processing (ICIP)}, 2017, pp. 3645–3649.
\bibitem{r21} Y. Zhang \emph{et al.}, ``ByteTrack: Multi-object tracking by associating every detection box,'' in \emph{Proc. European Conf. Computer Vision (ECCV)}, 2022.
\bibitem{r22} J. Cao, J. Pang, X. Weng, R. Khirodkar, and K. Kitani, ``Observation-centric SORT: Rethinking SORT for robust multi-object tracking,'' in \emph{Proc. IEEE/CVF Conf. Computer Vision and Pattern Recognition (CVPR)}, 2023, pp. 9686–9696.
\bibitem{strongsort} Y. Du, Z. Zhao, Y. Song, Y. Zhao, F. Su, T. Gong, and H. Meng, ``StrongSORT: Make DeepSORT great again,'' \emph{IEEE Trans. Multimedia}, vol. 25, 2023.
\bibitem{amodal} Y. Zhu, Y. Tian, D. Metaxas, and P. Dollár, ``Semantic amodal segmentation,'' in \emph{Proc. IEEE Conf. Computer Vision and Pattern Recognition (CVPR)}, 2017.
\bibitem{r23} A. Rudenko, L. Palmieri, M. Herman, K. M. Kitani, D. M. Gavrila, and K. O. Arras, ``Human motion trajectory prediction: A survey,'' \emph{Int. J. Robotics Research}, vol. 39, no. 8, pp. 895–935, 2020.
\bibitem{r24} A. Alahi, K. Goel, V. Ramanathan, A. Robicquet, L. Fei-Fei, and S. Savarese, ``Social LSTM: Human trajectory prediction in crowded spaces,'' in \emph{Proc. IEEE Conf. Computer Vision and Pattern Recognition (CVPR)}, 2016, pp. 961–971.
\bibitem{r25} T. Salzmann, B. Ivanovic, P. Chakravarty, and M. Pavone, ``Trajectron++: Dynamically-feasible trajectory forecasting with heterogeneous data,'' in \emph{Proc. European Conf. Computer Vision (ECCV)}, 2020.
\bibitem{r26} K. M. Kitani, B. D. Ziebart, J. A. Bagnell, and M. Hebert, ``Activity forecasting,'' in \emph{Proc. European Conf. Computer Vision (ECCV)}, 2012.
\bibitem{r27} D. Simon, ``Kalman filtering with state constraints: A survey of linear and nonlinear algorithms,'' \emph{IET Control Theory \& Applications}, vol. 4, no. 8, pp. 1303–1318, 2010.
\bibitem{r28} T. Nägeli, J. Alonso-Mora, A. Domahidi, D. Rus, and O. Hilliges, ``Real-time motion planning for aerial videography with dynamic obstacle avoidance and viewpoint optimization,'' \emph{IEEE Robotics and Automation Letters}, vol. 2, no. 3, pp. 1696–1703, 2017.
\bibitem{r29} B. Penin, P. Robuffo Giordano, and F. Chaumette, ``Vision-based reactive planning for aggressive target tracking while avoiding collisions and occlusions,'' \emph{IEEE Robotics and Automation Letters}, vol. 3, no. 4, pp. 3725–3732, 2018.
\bibitem{r30} R. Bonatti, W. Wang, C. Ho, A. Ahuja, M. Gschwindt, E. Camci, E. Kayacan, S. Choudhury, and S. Scherer, ``Autonomous aerial cinematography in unstructured environments with learned artistic decision-making,'' \emph{J. Field Robotics}, vol. 37, no. 4, pp. 606–641, 2020.
\bibitem{r31} Z. Han, R. Zhang, N. Pan, C. Xu, and F. Gao, ``Fast-Tracker: A robust aerial system for tracking agile target in cluttered environments,'' in \emph{Proc. IEEE Int. Conf. Robotics and Automation (ICRA)}, 2021, pp. 328–334.
\bibitem{r32} J. Ji, N. Pan, C. Xu, and F. Gao, ``Elastic Tracker: A spatio-temporal trajectory planner for flexible aerial tracking,'' in \emph{Proc. IEEE Int. Conf. Robotics and Automation (ICRA)}, 2022.
\bibitem{r33} H. Oleynikova, Z. Taylor, M. Fehr, R. Siegwart, and J. Nieto, ``Voxblox: Incremental 3D Euclidean signed distance fields for on-board MAV planning,'' in \emph{Proc. IEEE/RSJ Int. Conf. Intelligent Robots and Systems (IROS)}, 2017.
\bibitem{r34} A. D. Ames, S. Coogan, M. Egerstedt, G. Notomista, K. Sreenath, and P. Tabuada, ``Control barrier functions: Theory and applications,'' in \emph{Proc. European Control Conf. (ECC)}, 2019, pp. 3420–3431.
\end{thebibliography}

\end{document}
